\documentclass[11pt,a4paper]{article}
\usepackage[utf8]{inputenc}
\usepackage[T1]{fontenc}
\usepackage{amsmath,amssymb,amsthm}
\usepackage{graphicx}
\usepackage{hyperref}
\usepackage{booktabs}
\usepackage{geometry}
\usepackage{pgfplots}
\pgfplotsset{compat=1.18}
\geometry{margin=1in}

\newtheorem{theorem}{Theorem}[section]
\newtheorem{lemma}[theorem]{Lemma}
\newtheorem{proposition}[theorem]{Proposition}
\newtheorem{corollary}[theorem]{Corollary}
\newtheorem{definition}[theorem]{Definition}
\newtheorem{assumption}[theorem]{Assumption}
\newtheorem{remark}[theorem]{Remark}

\title{\textbf{Security Analysis of the Rand Protocol}\\[0.5em]
\large Post-Quantum Cryptographic Foundations, Attack Resistance,\\and Formal Security Guarantees}

\author{Dendi Suhubdy\\
Bitwyre Research\\
\texttt{dendi@bitwyre.com}}

\date{December 21, 2025}

\begin{document}

\maketitle

\begin{abstract}
We present a security analysis of the Rand protocol, a shielded pool deployed as a set of programs on Solana. Because Rand operates no validator set of its own, this analysis differs in scope from that of a Layer 1: consensus safety and liveness are inherited from Solana and are treated as assumptions rather than results, and we state explicitly where Rand's guarantees are conditional on them.

Our results concern the privacy construction and its cryptographic foundations. We prove transaction unlinkability and nullifier unlinkability under the zero-knowledge property of Groth16 and the collision resistance of Poseidon, and we establish conditional double-spend resistance from nullifier account uniqueness given Solana finality. We give a formal account of Rand's asymmetric post-quantum posture: note confidentiality, commitment hiding, and spend authorisation are post-quantum under lattice and QROM assumptions, while proof soundness rests on BN254 pairings and is therefore classical. We prove that this asymmetry is the correct allocation of a fixed budget, on the grounds that published ciphertexts are permanently exposed whereas an upgradeable verifier is not, and we state plainly the attack it leaves open.

We further enumerate leaks lying outside the proof system: public asset identifiers and fees, relayer-observable network metadata, anonymity set size, and note discovery. We regard understating these as a more serious defect in a privacy system than the leaks themselves.
\end{abstract}

\noindent\textbf{Keywords:} Post-Quantum Cryptography, Shielded Pool, Zero-Knowledge Proofs, Groth16, Lattice Cryptography, Solana, Metadata Privacy

\tableofcontents
\newpage

%==============================================================================
\section{Introduction}
%==============================================================================

The security of blockchain systems rests on cryptographic assumptions that face unprecedented challenges from advances in quantum computing. Shor's algorithm renders RSA and elliptic curve cryptography insecure against quantum adversaries, while Grover's algorithm reduces the effective security of symmetric primitives by half.

This paper establishes, for the Rand protocol:
\begin{enumerate}
    \item Post-quantum security of the confidentiality primitives under lattice and QROM assumptions
    \item A formal argument for treating soundness and confidentiality asymmetrically, and an honest statement of the residual quantum attack
    \item Privacy guarantees for the shielded transfer construction, together with an enumeration of leaks outside the proof system
    \item Conditional double-spend resistance, with the Solana assumptions it depends on made explicit
\end{enumerate}

\begin{remark}[Scope]
Rand runs as programs on Solana and operates no validator set. Consensus-level analysis, which occupied a substantial part of earlier versions of this document, is therefore not Rand's to provide: safety and liveness of ordering are Solana's, and appear below as assumptions. We consider it preferable to state this than to restate proofs about a consensus protocol Rand does not run.
\end{remark}

\subsection{Security Model}

\begin{definition}[Adversarial Models]
We consider adversaries of increasing capability:
\begin{enumerate}
    \item $\mathcal{A}_{\text{PPT}}$: Classical probabilistic polynomial-time adversary
    \item $\mathcal{A}_{\text{QPT}}$: Quantum polynomial-time adversary with quantum random oracle access
    \item $\mathcal{A}_{\text{Chain}}$: Adversary able to influence Solana's ordering, up to and including reorganisation of finalised history
    \item $\mathcal{A}_{\text{Economic}}$: Rational adversary maximizing profit
\end{enumerate}
\end{definition}

%==============================================================================
\section{Cryptographic Primitive Security}
%==============================================================================

\subsection{Hash Function Security}

We employ SHA-3 (Keccak-256) and BLAKE3 as collision-resistant hash functions.

\begin{definition}[Collision Resistance]
A hash function $H : \{0,1\}^* \to \{0,1\}^n$ is $(t, \epsilon)$-collision resistant if for all adversaries $\mathcal{A}$ running in time $t$:
\[
\Pr[(x, x') \leftarrow \mathcal{A} : x \neq x' \land H(x) = H(x')] \leq \epsilon
\]
\end{definition}

\begin{theorem}[Classical Collision Resistance]
SHA-3-256 is $(2^{128}, 2^{-128})$-collision resistant in the random oracle model.
\end{theorem}

\begin{theorem}[Quantum Collision Resistance]
SHA-3-256 achieves $(2^{85}, 2^{-85})$-collision resistance against quantum adversaries in the QROM.
\end{theorem}

\begin{proof}
The BHT algorithm finds collisions in $O(2^{n/3})$ quantum queries for an $n$-bit hash. For $n = 256$, this gives approximately $2^{85}$ queries.
\end{proof}

\subsection{Digital Signature Security: CRYSTALS-Dilithium}

We employ CRYSTALS-Dilithium, a lattice-based signature scheme based on Module-LWE and Module-SIS problems.

\begin{definition}[Module-LWE Problem]
Let $R_q = \mathbb{Z}_q[X]/(X^n + 1)$ be a polynomial ring. The Module-LWE problem with parameters $(n, k, q, \chi)$ is: given $(\mathbf{A}, \mathbf{b} = \mathbf{A}\mathbf{s} + \mathbf{e})$ where $\mathbf{A} \leftarrow R_q^{k \times k}$, $\mathbf{s} \leftarrow \chi^k$, $\mathbf{e} \leftarrow \chi^k$, distinguish this from uniform.
\end{definition}

\begin{assumption}[Hardness of Module-LWE]
For parameters $(n = 256, k = 4, q = 8380417, \sigma = 3.2)$, there exists no QPT adversary that solves Module-LWE with advantage greater than $2^{-128}$.
\end{assumption}

\begin{theorem}[Dilithium EUF-CMA Security]
Under the Module-LWE assumption, CRYSTALS-Dilithium is existentially unforgeable under chosen message attacks with security $2^{-128}$ against QPT adversaries.
\end{theorem}

\begin{table}[h]
\centering
\caption{Dilithium Security Levels}
\begin{tabular}{@{}lccc@{}}
\toprule
\textbf{Level} & \textbf{Classical} & \textbf{Quantum} & \textbf{Parameters} \\ \midrule
Dilithium2 & 128 bits & 128 bits & $k = 4, \ell = 4$ \\
Dilithium3 & 192 bits & 192 bits & $k = 6, \ell = 5$ \\
Dilithium5 & 256 bits & 256 bits & $k = 8, \ell = 7$ \\ \bottomrule
\end{tabular}
\end{table}

Rand uses Dilithium2 for transaction signatures.

%==============================================================================
\section{Post-Quantum Security Analysis}
%==============================================================================

\subsection{Quantum Adversary Model}

\begin{definition}[Quantum Polynomial-Time]
An adversary $\mathcal{A}_{\text{QPT}}$ has access to:
\begin{enumerate}
    \item A fault-tolerant quantum computer with $\text{poly}(\lambda)$ qubits
    \item Quantum superposition queries to random oracles (QROM)
    \item Classical communication with other protocol participants
\end{enumerate}
\end{definition}

\subsection{Shor's Algorithm Resistance}

\begin{theorem}[Resistance to Shor's Algorithm]
The Rand protocol is immune to Shor's algorithm because it uses no primitives based on:
\begin{enumerate}
    \item Integer factorization (RSA)
    \item Discrete logarithm in finite fields (DH, DSA)
    \item Discrete logarithm in elliptic curves (ECDSA, ECDH)
\end{enumerate}
\end{theorem}

\begin{proof}
We enumerate all cryptographic operations in Rand:
\begin{itemize}
    \item Signatures: CRYSTALS-Dilithium (lattice-based) - immune to Shor
    \item Key exchange: CRYSTALS-Kyber (lattice-based) - immune to Shor
    \item Hashing: SHA-3, BLAKE3 (symmetric) - Shor does not apply
    \item Note commitments: Poseidon (hash-based) - Shor does not apply
    \item ZK proof soundness: Groth16 over BN254 - \textbf{broken by Shor}, see Section \ref{sec:asym}
\end{itemize}
With one exception. Groth16 verification evaluates pairings on BN254, whose security rests on the elliptic-curve discrete logarithm problem, and Shor's algorithm solves it. Rand's proof soundness is therefore \emph{not} post-quantum. Section \ref{sec:asym} argues that this is the correct place to absorb the compromise, and Section \ref{sec:residual} states precisely what it permits an adversary to do.
\end{proof}

\subsection{Grover's Algorithm Impact}

\begin{theorem}[Grover Security Bounds]
Grover's algorithm provides quadratic speedup for unstructured search:
\end{theorem}

\begin{table}[h]
\centering
\caption{Grover's Impact on Primitives}
\begin{tabular}{@{}lccc@{}}
\toprule
\textbf{Primitive} & \textbf{Classical} & \textbf{Quantum} & \textbf{Our Choice} \\ \midrule
Hash preimage & $2^{256}$ & $2^{128}$ & SHA-3-256 \\
Hash collision & $2^{128}$ & $2^{85}$ & SHA-3-256 \\
AES key search & $2^{256}$ & $2^{128}$ & AES-256-GCM \\
Dilithium forgery & $2^{128}$ & $2^{128}$* & Level 2 \\ \bottomrule
\end{tabular}
\end{table}

*Lattice problems are not amenable to Grover speedup.

\subsection{Asymmetric Post-Quantum Posture}
\label{sec:asym}

Rand does not claim uniform post-quantum security. It claims post-quantum \emph{confidentiality} alongside classical \emph{soundness}, and this section argues that the split is principled rather than expedient.

The relevant distinction is temporal. Some security properties, once violated, are violated retroactively over all data ever published; others are re-established at each use and can be strengthened before the violation occurs.

\begin{definition}[Irrevocable and Renewable Properties]
A security property is \emph{irrevocable} if an adversary who obtains the relevant capability at time $t_1$ can violate it with respect to data published at any $t_0 < t_1$. It is \emph{renewable} if violation requires the capability at the moment of use, so that replacing the mechanism before $t_1$ preserves the property.
\end{definition}

\begin{theorem}[Correct Allocation of Post-Quantum Budget]
\label{thm:alloc}
Given a construction in which post-quantum security can be afforded for either confidentiality or soundness but not both, allocating it to confidentiality strictly dominates.
\end{theorem}

\begin{proof}
Confidentiality of published ciphertexts is irrevocable. A note ciphertext committed to the ledger at $t_0$ is recorded permanently by any observer; an adversary acquiring quantum capability at $t_1 > t_0$ decrypts it, and no action during $(t_0, t_1)$ prevents this, since publication cannot be undone. Allocating post-quantum security elsewhere therefore forfeits confidentiality of all history.

Soundness is renewable. It is enforced by a verifier executed afresh at each verification, and Solana programs are upgradeable, so the verifier may be replaced at any $t < t_1$. Proofs verified before the replacement were checked under an assumption then unbroken; proofs after it are checked under the new one. Post-quantum soundness is thus obtainable later at no retroactive cost.

Allocating to confidentiality secures an irrevocable property immediately and defers a renewable one; allocating to soundness secures a renewable property that could have been deferred and permanently forfeits an irrevocable one. The former dominates.
\end{proof}

\begin{table}[h]
\centering
\caption{Post-Quantum Status by Property}
\begin{tabular}{@{}llll@{}}
\toprule
\textbf{Property} & \textbf{Primitive} & \textbf{Class} & \textbf{PQ} \\ \midrule
Note confidentiality & ML-KEM (Kyber768) & irrevocable & yes \\
Commitment hiding & Poseidon & irrevocable & yes \\
Spend authorisation & ML-DSA (Dilithium2) & irrevocable & yes \\
Proof soundness & Groth16 / BN254 & renewable & \textbf{no} \\ \bottomrule
\end{tabular}
\end{table}

\subsection{Residual Quantum Attack}
\label{sec:residual}

We state the consequence of Theorem \ref{thm:alloc} without softening it.

\begin{proposition}[Residual Attack Surface]
An adversary possessing a cryptographically relevant quantum computer, acting before the verifier is upgraded, can compute discrete logarithms on BN254, forge a Groth16 proof for a false statement, and thereby withdraw value from the shielded pool without owning corresponding notes. Such an adversary cannot decrypt historical note ciphertexts, recover spend keys from published commitments, or deanonymise past transfers.
\end{proposition}

The exposure is therefore to theft of pool balances at the time of attack, not to retrospective loss of privacy. We regard this as the correct risk to carry, since the former is bounded by pool size and detectable through supply auditing, while the latter would be unbounded, silent, and irreversible.

\begin{remark}[Migration trigger]
The verifier should be replaced once a succinct post-quantum argument system exists whose proofs fit a Solana transaction and whose verification fits the compute budget. Neither lattice-based nor hash-based succinct arguments met both conditions at the time of writing; the binding constraint is the 1232-byte transaction limit rather than proving cost.
\end{remark}

\subsection{Overall Post-Quantum Security}

\begin{theorem}[Post-Quantum Confidentiality]
Rand achieves 128-bit \emph{confidentiality} against quantum polynomial-time adversaries under:
\begin{enumerate}
    \item Module-LWE with parameters $(n = 256, k = 4, q = 8380417)$
    \item Module-SIS with the same parameters
    \item SHA-3, BLAKE3, and Poseidon modelled as quantum random oracles
\end{enumerate}
\end{theorem}

\begin{remark}
This theorem is deliberately narrower than the corresponding claim in earlier versions of this document, which asserted 128-bit post-quantum security for the protocol as a whole. That claim is not supported once soundness rests on BN254, and it has been withdrawn rather than qualified.
\end{remark}

%==============================================================================
\section{Inherited Consensus Assumptions}
%==============================================================================

Rand runs no consensus protocol. This section states what is assumed of Solana and what fails if the assumption fails.

\begin{assumption}[Solana Finality]
\label{ass:finality}
\label{asm:finality}
A transaction finalised by Solana is never reordered or reverted.
\end{assumption}

\begin{assumption}[Solana Liveness]
\label{ass:liveness}
Submitted transactions paying sufficient priority fee are eventually included.
\end{assumption}

\begin{theorem}[Conditional Double-Spend Resistance]
\label{thm:nodouble}
Under Assumption \ref{asm:finality}, no note in the shielded pool can be spent twice.
\end{theorem}

\begin{proof}
A note determines a unique nullifier $\mathsf{nf} = \mathsf{Poseidon}(\mathsf{nk} \| \rho)$, and the circuit constrains its derivation, so a spend of a given note cannot present any other value. The pool program creates a program-derived account seeded by $\mathsf{nf}$; the runtime fails account creation if the account already exists. By Assumption \ref{asm:finality} the two spend attempts are totally ordered; the earlier creates the account and succeeds, the later attempts creation of an existing account and aborts. Hence at most one spend of any note is finalised.
\end{proof}

\begin{remark}[Failure mode]
Theorem \ref{thm:nodouble} is conditional, and the condition is load-bearing. Were finalised Solana history reorganised, nullifier account creations could be reordered or dropped and the argument would not hold. Rand has no independent defence: its safety ceiling is Solana's. This is the direct cost of not operating a validator set, and it should be weighed against the corresponding benefits, namely native asset access and no separate security budget.
\end{remark}

\begin{remark}[Sybil resistance]
Rand requires no Sybil resistance mechanism. It has no validator set to join, no stake to acquire, and no leader election to influence. Relaying is permissionless by design and a relayer cannot alter the operation it submits, since every public input is bound by the proof; the worst behaviour available to a relayer is to decline, which leaves the operation available to others.
\end{remark}

%==============================================================================
\section{Privacy Security}
%==============================================================================

\subsection{Transaction Privacy Model}

\begin{definition}[Transaction Unlinkability]
A private transaction scheme satisfies unlinkability if for any PPT adversary $\mathcal{A}$:
\[
\Pr[\mathcal{A}(\text{tx}_0, \text{tx}_1) = 1] - \frac{1}{2} \leq \text{negl}(\lambda)
\]
where $\text{tx}_0$ and $\text{tx}_1$ are transactions with different hidden senders/recipients/amounts.
\end{definition}

\begin{theorem}[Groth16 Zero-Knowledge]
Groth16 satisfies computational zero-knowledge: there exists a PPT simulator $\mathcal{S}$ such that for all $(x, w) \in R$:
\[
\{\text{Prove}(x, w)\} \approx_c \{\mathcal{S}(x)\}
\]
\end{theorem}

\begin{theorem}[Nullifier Unlinkability]
Nullifiers $\mathsf{nf} = \mathsf{Poseidon}(\mathsf{nk} \| \rho)$ are unlinkable to their corresponding commitments under the PRF security of Poseidon.
\end{theorem}

\begin{theorem}[Transaction Privacy]
A shielded transfer reveals only:
\begin{enumerate}
    \item The Merkle root, hence which anonymity set is used
    \item Nullifiers, hence that some notes are spent
    \item Output commitments, hence that some notes are created
    \item The asset identifier and the SHRUGG fee
\end{enumerate}
and nothing about sender identity, recipient identity, or amounts.
\end{theorem}

\subsection{Leaks Outside the Proof System}

Theorem statements about the circuit do not bound what an observer learns in practice. We enumerate the gaps.

\begin{table}[h]
\centering
\caption{Metadata Leaks and Mitigations}
\begin{tabular}{@{}llp{4.6cm}@{}}
\toprule
\textbf{Leak} & \textbf{Observable by} & \textbf{Mitigation} \\ \midrule
Asset identifier & anyone & none; inherent to a multi-asset pool \\
Fee amount & anyone & fee quantisation to reduce distinguishability \\
Network origin & relayer & anonymising transport \\
Anonymity set size & anyone & none; improves only with adoption \\
Note discovery scan & indexer, if delegated & self-scan; fuzzy message detection \\
\textbf{Full witness} & \textbf{prover, if delegated} & \textbf{self-prove; blinded delegation (future)} \\
Deposit and withdrawal & anyone & timing separation, amount quantisation \\ \bottomrule
\end{tabular}
\end{table}

\begin{theorem}[Witness Disclosure Under Delegation]
\label{thm:witness-sec}
A prover supplied with the witness for a transfer learns the sender's spend authority, the recipient, and the amount, irrespective of the zero-knowledge property of the resulting argument.
\end{theorem}

\begin{proof}
Zero-knowledge bounds what the argument $\pi$ reveals to a verifier holding only the statement $x$; it constrains nothing about a party holding the witness $w$ directly. Since $w$ comprises the spent notes' openings, the authorising key material, and the output values and recipients, a prover evaluating the circuit on $w$ reads exactly these quantities.
\end{proof}

Three leaks deserve emphasis. First, transfers are formally private within an anonymity set whose size the protocol cannot guarantee; a transfer in a pool with few active participants is identifiable in practice irrespective of the cryptography, and early deployments are therefore weakest precisely when users may assume otherwise. Second, deposits and withdrawals are transparent at the pool boundary, so correlating an entry with a later exit of similar magnitude and timing remains the most effective attack against the system. Third, delegated proving discloses everything the circuit conceals, to a party of the user's choosing --- a strictly different threat model from the others in this table, since it is opt-in and identifiable rather than ambient, but the strongest disclosure available to any single party in the system. None is a defect in the proof system, and none is fixed by strengthening it.

%==============================================================================
\section{Attack Resistance Analysis}
%==============================================================================

\subsection{Consensus-Level Attacks}

Attacks on ordering and Sybil resistance apply to Solana, not to Rand, and Rand neither defends against them nor claims to. Their treatment belongs to analyses of Solana's consensus.

\begin{table}[h]
\centering
\caption{Consensus-Level Attacks Under Inheritance}
\begin{tabular}{@{}lll@{}}
\toprule
\textbf{Attack} & \textbf{Applies to Rand?} & \textbf{Consequence if successful} \\ \midrule
51\% / stake majority & no, to Solana & Assumption \ref{asm:finality} fails; double-spend possible \\
Long-range & no, to Solana & as above \\
Nothing-at-stake & no, to Solana & as above \\
Grinding & no, to Solana & leader prediction; no Rand-specific effect \\
\bottomrule
\end{tabular}
\end{table}

Every entry reduces to the same statement: Rand's safety is conditional on Solana's, and a successful consensus attack on Solana invalidates Theorem \ref{thm:nodouble}. Restating BFT threshold arguments here would imply Rand enforces them, which it does not.

\subsection{Network-Level Attacks}

\subsubsection{Eclipse Attack}

\begin{theorem}[Eclipse Attack Mitigation]
With $k$ outbound connections to diverse subnets, an adversary must control nodes in $k$ different /16 networks:
\[
\Pr[\text{Eclipse}] \leq \left(\frac{m}{N}\right)^k
\]
where $m$ is adversary-controlled nodes and $N$ is total network size.
\end{theorem}

\subsubsection{Sybil Attack}

\begin{theorem}[Inherited Sybil Resistance]
Rand implements no Sybil-resistance mechanism. Consensus participation is Solana's, and its stake-weighted leader schedule provides
\[
\text{Power}(k \text{ identities}) = k \cdot \frac{S}{k} = S = \text{Power}(1 \text{ identity})
\]
so splitting stake provides no benefit. Rand's own roles --- relaying and proving --- are deliberately Sybil-indifferent: both are permissionless and unbonded, and operating $k$ identities confers no advantage over operating one, since the nullifier set admits exactly one submission per batch regardless of who submits it.
\end{theorem}

\begin{remark}[Correction to earlier revisions]
Previous versions of this document stated Sybil resistance as a property Rand established through stake-weighted consensus of its own. That was an artefact of the superseded Layer 1 design. Rand operates no validator set and holds no stake; the property is inherited, and it is bounded by Assumptions \ref{ass:finality} and \ref{ass:liveness} rather than by anything Rand enforces.
\end{remark}

\subsection{Transaction-Level Attacks}

\subsubsection{Front-Running / MEV}

\begin{theorem}[Front-Running Mitigation]
A shielded transfer's contents are concealed until after inclusion, so an adversary observing the mempool cannot condition an ordering strategy on the transfer's sender, recipient, or amount.
\end{theorem}

\begin{remark}[The claim is conditional, not absolute]
Earlier revisions asserted that private transactions \emph{eliminate} front-running. That is too strong. What is concealed is the transfer's contents, not its existence, its asset identifier, its fee, or its timing --- and an adversary may still order or censor on those. Where proving is delegated, the prover additionally holds the witness by Theorem \ref{thm:witness-sec} before submission, and is therefore in precisely the position a front-runner would wish to occupy. Rand does not order its own transactions and cannot exclude ordering strategies available to Solana's leader. The honest statement is mitigation, bounded by what remains observable.
\end{remark}

\subsubsection{Double-Spend Attack}

\begin{theorem}[Double-Spend Prevention]
The combination of Solana's finality, per Assumption \ref{ass:finality}, and nullifier uniqueness prevents double-spending. Nullifier accounts are created on spend, and account creation is idempotent within Solana's runtime, so a second spend of the same note fails at account creation rather than at proof verification.
\end{theorem}

\subsection{Attack Resistance Summary}

\begin{table}[h]
\centering
\caption{Attack Resistance Summary}
\begin{tabular}{@{}lccc@{}}
\toprule
\textbf{Attack} & \textbf{Classical} & \textbf{Quantum} & \textbf{Mechanism} \\ \midrule
51\% Attack & 66.7\% threshold & Same & BFT \\
Long-Range & Checkpointing & Same & Weak subjectivity \\
Nothing-at-Stake & 33\% slashing & Same & Economic penalty \\
Grinding & VRF & Same & Unpredictable \\
Eclipse & Peer diversity & Same & Network topology \\
Sybil & Stake-weighted & Same & Economic cost \\
Front-Running & Private txs & Same & ZK proofs \\
Replay & Nullifiers & Same & Uniqueness \\
Double-Spend & BFT + nullifiers & Same & Finality \\
Quantum (Shor) & N/A & Lattice-based & PQ crypto \\
Quantum (Grover) & N/A & 256-bit primitives & Larger keys \\ \bottomrule
\end{tabular}
\end{table}

%==============================================================================
\section{Economic Security}
%==============================================================================

Rand has no stake, so the customary analysis, in which attack cost is bounded below by the price of acquiring a stake fraction, does not apply. Two economic questions take its place.

\subsection{Relayer Market Liveness}

Shielded operations reach the chain only if some relayer finds submission profitable. With in-circuit fee $f$ in SHRUGG and cost $c$ in SOL, a relayer submits when $f \cdot p_{\text{SHRUGG}} > c \cdot p_{\text{SOL}}$.

\begin{proposition}[Liveness Under Rational Relayers]
If at least one rational relayer observes the queue and users set $f$ such that $f \cdot p_{\text{SHRUGG}}$ exceeds prevailing costs, every operation is eventually submitted.
\end{proposition}

The dependence runs on the price ratio. A sharp fall in SHRUGG against SOL, or a congestion spike in priority fees, can render outstanding operations unprofitable, stalling them until users re-price. This is a liveness risk without a safety component: a stalled operation is not lost, since its notes remain unspent and it may be resubmitted at a higher fee. Users should nonetheless expect degraded service under adverse conditions rather than a hard guarantee.

\subsection{Pool Balance as Attack Incentive}

Absent staking, the economically relevant target is the pool itself. Its balance is public, since deposits and withdrawals are transparent, so the reward for a successful soundness break is known to any attacker in advance.

\begin{remark}
This is a genuine and structural weakness. A shielded pool is a concentrated, publicly quantified honeypot whose exploitation depends on a single soundness assumption and one program's correctness. It argues for a conservative deposit cap during early operation, for supply auditing sufficient to detect unbacked withdrawal quickly, and for a bounded rather than open-ended migration path away from BN254. We prefer to state this plainly rather than present a stake-based cost bound that does not apply.
\end{remark}

%==============================================================================
\section{Concrete Parameters}
%==============================================================================

\begin{table}[h]
\centering
\caption{Recommended Parameters for 128-bit Security}
\begin{tabular}{@{}lll@{}}
\toprule
\textbf{Component} & \textbf{Parameter} & \textbf{Security Level} \\ \midrule
\multicolumn{3}{c}{\textit{Cryptographic Primitives}} \\
Hash function & SHA-3-256 & 128-bit classical, 85-bit quantum \\
Signatures & Dilithium2 & 128-bit quantum \\
Key encapsulation & Kyber768 & 192-bit quantum \\
Symmetric encryption & AES-256-GCM & 128-bit quantum \\
\midrule
\multicolumn{3}{c}{\textit{Proof System Parameters}} \\
Argument system & Groth16 over BN254 & \textbf{classical} soundness \\
In-circuit hash & Poseidon & 128-bit hiding \\
Proof size & 256 bytes & fits one Solana transaction \\
Setup & per-circuit MPC ceremony & secure if one party honest \\
\midrule
\multicolumn{3}{c}{\textit{Settlement Parameters (inherited)}} \\
Ordering & Solana Proof of History & assumption, not result \\
Finality & Solana schedule & Assumption \ref{asm:finality} \\
Recent-root window & 64 & concurrency tolerance \\
\midrule
\multicolumn{3}{c}{\textit{Privacy Parameters}} \\
Merkle tree depth & 32 & $2^{32}$ notes capacity \\
Nullifier size & 256 bits & 128-bit collision resistance \\
Note commitment & 256 bits & 128-bit hiding \\
\bottomrule
\end{tabular}
\end{table}

%==============================================================================
\section{Conclusion}
%==============================================================================

We have analysed the security of the Rand shielded pool, establishing:

\begin{enumerate}
    \item 128-bit post-quantum \emph{confidentiality} under lattice and QROM assumptions
    \item That allocating a limited post-quantum budget to confidentiality rather than soundness strictly dominates (Theorem \ref{thm:alloc}), together with an explicit statement of the attack this leaves open
    \item Transaction and nullifier unlinkability under Groth16 zero-knowledge and Poseidon collision resistance
    \item Conditional double-spend resistance (Theorem \ref{thm:nodouble}), with its dependence on Solana finality made explicit
\end{enumerate}

We make no claim to the strongest security guarantees available. Rand's safety is bounded above by Solana's, its proof soundness is classical pending a practical post-quantum succinct argument, and its practical privacy is bounded by anonymity set size and by boundary correlation between deposits and withdrawals rather than by the cryptography. These are the honest limits of the construction, and a privacy system whose stated guarantees exceed its real ones is more dangerous to its users than one with modest and accurate claims.

Future work includes the trusted setup ceremony, formal verification of the circuit, mechanised proofs of the pool program's state transitions, and monitoring for a post-quantum argument system whose verification fits a Solana transaction.

\bibliographystyle{plain}
\begin{thebibliography}{9}
\bibitem{dilithium} Ducas et al., ``CRYSTALS-Dilithium,'' TCHES 2018.
\bibitem{kyber} Bos et al., ``CRYSTALS-Kyber,'' Euro S\&P 2018.
\bibitem{stark} Ben-Sasson et al., ``Scalable, transparent, and post-quantum secure computational integrity,'' CRYPTO 2018.
\bibitem{groth16} J. Groth, ``On the size of pairing-based non-interactive arguments,'' EUROCRYPT 2016.
\bibitem{poseidon} L. Grassi et al., ``Poseidon: A new hash function for zero-knowledge proof systems,'' USENIX Security 2021.
\bibitem{fmd} G. Beck, J. Len, I. Miers, and M. Green, ``Fuzzy message detection,'' ACM CCS 2021.
\bibitem{grover} Grover, ``A fast quantum mechanical algorithm for database search,'' STOC 1996.
\bibitem{shor} Shor, ``Algorithms for quantum computation,'' FOCS 1994.
\end{thebibliography}

\end{document}
